% Section 8: Conclusion (about 0.3 page). Plan: PAPER_PLAN 4 / §8.
% A check-in-level representation makes the category task more learnable; one model outperforms the
% dedicated region model where the number of regions is large and matches it where it is small, and
% on category outperforms at FL and holds within a fifth of a point elsewhere. One model with a
% shared semantic context and a private spatial path anticipates both what and where.
\section{Conclusion}
\label{sec:conclusion}

We tested whether one model can predict both the category and region of a user's next visit. A check-in-level
representation, one vector per visit rather than one per place, makes the next-category task far more learnable than the standard
place-level representation does. On that representation, a single multi-task model reads both answers from one saved checkpoint. On region, it outperforms the dedicated model at Texas and California, the two datasets with the largest region vocabularies, and remains non-inferior (TOST, $\pm2$~pp) at the other four. On category, it outperforms the dedicated model at Florida and stays within a fifth of a point of it elsewhere. Joint training therefore pays where the region task is hardest and costs nothing measurable where it is not. The joint model also remains above every external baseline in Table~\ref{tab:results} on all six datasets, by at least $3.5$ Acc@10 points over the strongest region reference (HMT-GRN, STAN, or ReHDM; the first two on our folds, ReHDM under its own protocol) and by at least $3.0$ macro-F1 points over POI-RGNN on category. One model, with a shared semantic context and a private spatial path, anticipates both what a user will do next and where, at the accuracy two dedicated models reach and at one forward pass instead of two.
